\subsection{Regular values and transversality}
    It is known that for the regular value theorem, and more generally statements about transversality, to still hold for manifolds with boundary one needs to impose additional conditions on the boundary behaviour of a smooth function $f \colon M \to N$. Here is another lemma, that fits into the same category. We call $q \in N$ a \textbf{boundary value} of $f$ if $q \in f(\partial M)$.

    \begin{lemma}
        \label{lem:regular_sublevel_sets}
        Let $M$ be a manifold with boundary and $f \colon M \to \R$ a smooth function. Furthermore, let $a \in \R$ be a regular value of $f$, that is not a boundary value. Then $M_a := f^{-1}(-\infty,a]$ is a submanifold with boundary 
        \[
            \partial M_a = f^{-1}(a) \cup (M_a \cap \partial M).
        \]
    \end{lemma}
    \begin{proof}[Proof sketch]
        For $p \in f^{-1}(a)$ we have $p \in \mathrm{int}(M)$ by assumption. Since $\mathrm{int}(M)$ is a boundaryless manifold (and an open subset of $M$) the local submersion theorem holds, i.e. $f$ looks like the orthogonal projection $\R^n \to \R$ around $p$. Then the proof goes essentially the same as in the boundaryless case.
    \end{proof}

    \begin{counterexample}
        \label{ex:boundary_values}
        To see that Lemma \ref{lem:regular_sublevel_sets} fails without the assumption that $a$ is not a boundary value consider $M := S^1 \times [0,1]$ and $f \colon M \to \R, (z,t) \mapsto t$. Then $f$ is everywhere regular but $f^{-1}(-\infty,0]$ does not satisfy the statements in Lemma \ref{lem:regular_sublevel_sets}. 
    \end{counterexample}

    Let $M,N$ be manifolds with boundary and $f \colon M\to N$ smooth. We define the \textbf{boundary map} $\partial f$ of $f$ by 
    \[
    \partial f := f|_{\partial M}.
    \]
    We now generalize three theorems from \cite{guillemin2010differential} about transversal maps by allowing the codomain to have boundary as well. 

    \begin{theorem}[Transversality Theorem]
    \label{thm:transversality}
    Let $M,N$ be manifolds with boundary, $f \colon M \to N$ smooth and $X \subseteq N$ a submanifold without boundary. If 
    \[
    f \pitchfork X \text{ and } \partial f \pitchfork X
    \]
    then $X^\prime := f^{-1}(X)$ is a submanifold with boundary 
    \[
    \partial X^\prime = X^\prime \cap \partial M
    \]
    and 
    \[
    \mathrm{codim}(X^\prime) = \mathrm{codim}(X).
    \]
    \end{theorem}
    \begin{proof}
        The case $\partial N = \emptyset$ is shown in \cite{guillemin2010differential} (p. 60). If $\partial N \neq \emptyset$ consider the inclusion
        \[
        \iota \colon N \to \mathrm{D}(N),
        \]
        where $\mathrm{D}(N)$ denotes the double of $N$ and define $F := \iota \circ f$. Then $Y := \iota(X)$ is a submanifold of $\mathrm{D}(N)$ and $X^\prime = F^{-1}(Y)$. Furthermore, since $\iota$ is an embedding (and $\mathrm{dim}(N) = \mathrm{dim}(\mathrm{D}(N))$) it follows that 
        \[
        F \pitchfork Y \text{ and } \partial F \pitchfork Y,
        \] 
        such that the statement follows from the boundaryless case.
    \end{proof}

       We can also make a statement about the tangent space of $X^\prime$ in the above theorem. To this end let $f \colon M \to N$ and $X \subseteq N, X^\prime \subseteq M$ be as above. Then we can also characterize transversality $f \pitchfork X$ as follows: Choose a metric on $N$ and for $x \in X^\prime, y = f(x)$ define the \textbf{vertical differential}
    \[
    D^V_xf \colon T_xM \to N_yX
    \]
    as the composition 
    \[
    T_xM \xrightarrow{d_xf} T_yN = T_yX \oplus N_yX \xrightarrow{\mathrm{proj.}} N_yX.
    \]
    Then $f \pitchfork X$ if and only if $D^V_xf$ is surjective for all $x \in X^\prime$. If this is the case we moreover have (at least if $\partial M = \emptyset$, otherwise we also require $\partial f \pitchfork X$)
    \[
    \mathrm{codim}(\mathrm{ker}D^V_xf) = \mathrm{dim}(\mathrm{im}(D^V_xf)) = \mathrm{codim}(X) = \mathrm{codim}(X^\prime)
    \]
    and thus (since $T_xX^\prime \subseteq \mathrm{ker}D^V_xf$) we get 
    \[
    T_xX^\prime = \mathrm{ker}D^V_xf.
    \]
    This generalizes the formula for the tangent space of the preimage of a regular value. It should also be noted, that the Riemannian metric on $N$ was only chosen, to define a complement of $T_yX$ in a consistent manner, such that the vertical differential $D^V_xf$ depends smoothly on $x$.

    \begin{theorem}[Parametric Transversality Theorem]
        \label{thm:parametric_transversality}
        Let $M,N$ be manifolds with boundary and $X \subseteq N$ a submanifold without boundary. Furthermore, let $S$ be a manifold without boundary and $F \colon M \times S \to N$ a smooth map with 
        \[
            F \pitchfork X \text{ and } \partial F \pitchfork X.
        \]
        Then for almost every $s \in S$ the map $f_s \colon M \to N$ defined by $f_s(p) := F(p,s)$ satisfies
        \[
            f_s \pitchfork X \text{ and } \partial f_s \pitchfork X.
        \]
    \end{theorem}
    \begin{proof}
        By Theorem \ref{thm:transversality} we know that 
        \[
        X^\prime := F^{-1}(X)
        \]
        is a submanifold (with boundary) of $M \times S$. Consider the projection $\pi \colon M \times S \to S$. We can then show, that for a regular value $s \in S$ of the map $\pi|_{X^\prime}$ the maps $f_s$ and $\partial f_s$ are transversal to $X$, which is just some linear algebra. The proof is exactly the same as in the case $\partial N = \emptyset$, see \cite{guillemin2010differential} (p. 68). Then the statement follows from Sard`s Theorem.
    \end{proof}

    \begin{theorem}[Transversality Homotopy Theorem]
        \label{thm:transversality_homotopy}
        Let $M,N$ be manifolds with boundary, $f \colon M \to N$ smooth and $X \subseteq N$ a submanifold without boundary. Then $f$ is smoothly homotopic to a map $g \colon M \to N$ with
        \[
            g \pitchfork X \text{ and } \partial g \pitchfork X.
        \]
    \end{theorem}
    \begin{proof}
        First we may assume that $\mathrm{im}(f)\subseteq \mathrm{int}(N)$, since $f$ is always smoothly homotopic to such a map (the inclusion $\mathrm{int}(N) \hookrightarrow N$ is a smooth homotopy equivalence). Define
        \[
        Y := X \cap \mathrm{int}(N)
        \]
        and 
        \[
        f^\prime \colon M \to \mathrm{int}(N), \; p \mapsto f(p).
        \]
        Using the Transversality Homotopy Theorem for the case $\partial N = \emptyset$ (see \cite{guillemin2010differential} p. 70) we get that $f^\prime$ is smoothly homotopic to a map $g^\prime \colon M \to \mathrm{int}(N)$, such that 
        \[
            g^\prime \pitchfork Y \text{ and } \partial g^\prime \pitchfork Y.
        \]
        Now define $g := \iota \circ f^\prime$, where $\iota \colon \mathrm{int}(N) \hookrightarrow N$ is the inclusion. Since $g$ does not intersect $X$ at boundary points, i.e. $\mathrm{im}(g) \cap X \subseteq \mathrm{int}(N)$, the transversality properties of $g^\prime$ directly imply
        \[
            g \pitchfork X \text{ and } \partial g \pitchfork X.
        \]
        
        %Now consider the inclusion $\iota \colon N \to \mathrm{D}(N)$ of $N$ into its double and define $g := \iota \circ f$ and $Y := \iota(X)$. From here on we mimic the proof of Theorem 6.36 in \cite{lee2011smooth}. By the Whitney embedding theorem we may assume that $\mathrm{D}(N)$ is a submanifold of some $\R^n$. Choose a tubular neighbourhood $V$ of $\mathrm{N}$ together with a retraction $r \colon V \to \mathrm{D}(N)$, which is a submersion. Now define $$\delta \colon \mathrm{D}(N) \to \R_{>0}, \; x \mapsto \mathrm{sup}\{\tau >0 \mid B_\tau(x) \subseteq V\}.$$
        % It is straightforward to show that $\delta$ is continuous. Using the Whitney approximation theorem for functions we can choose a smooth function $\epsilon \colon \mathrm{D}(N) \to \R_{>0}$ with 
        % \[
        % 0 < \epsilon(x) < \delta(x) \text{ for all } x \in \mathrm{D}(N). 
        % \]
        % Now define 
        % \[
        % G^\prime \colon M \times B_1(0) \to V, (p,v) \mapsto g(p) + \epsilon(g(p))v. 
        % \]
        % and $G := r \circ G^\prime$. For $p \in M$ consider the inclusion
        % \[
        % h_p \colon B_1(0) \to M \times B_1(0), \; v \mapsto (p,v).
        % \]
        % Then $G^\prime \circ h_p$ is clearly a submersion. If $p \in \partial M$ then $\partial G^\prime \circ h_p$ is also a submerison. In particular $D_{(p,v)} G^\prime$ and $D_{(p,v)} \partial G^\prime$ are surjective, thus $G^\prime$ and $\partial G^\prime$ are submersions. This implies that $G$ and $\partial G$ are submerisons (since $r$ is also a submerison). In particular 
        % \[
        % G \pitchfork Y \text{ and } \partial G \pitchfork Y.
        % \]
        % Applying Theorem \ref{thm:parametric_transversality} we get that for almost every $v \in B_1(0)$:
        % \[
        %     g_v \pitchfork Y \text{ and } \partial g_v \pitchfork Y.
        % \]
        % Since $g_0 = g = \iota \circ f$ and $\mathrm{im}(f) \subseteq \mathrm{int}(N)$ we get that $\mathrm{im}(g_0) \subseteq \iota(\mathrm{int}(N)) =: U$, which is an open subset of $\mathrm{D}(N)$. Then
        % \[
        % G^{-1}(U) \subseteq M \times B_1(0)
        % \]
        % is open and $M \times {0} \subseteq U$.
    \end{proof}

    In the above theorems the submanifold $X$ which we pulled back by $f$ was always assumed to be without boundary. Here is a version, which allows for $X$ to have boundary as well.

    \begin{lemma}
        \label{on_mnf_with_boundary:lem:preimage_of_submersion}
            Let $M,N$ be smooth manifolds where $M$ has no boundary and let $f \colon M \to N$ be a submersion. Given a submanifold $X \subseteq N$ with boundary the preimage $f^{-1}(X) \subseteq M$ is a submanifold with boundary 
            \[
            \partial f^{-1}(X) = f^{-1}(\partial X)
            \]
            and $\mathrm{codim}(f^{-1}(X)) = \mathrm{codim}(X)$. 
        \end{lemma}
        \begin{proof}
            First we may also assume $\partial N = \emptyset$, otherwise embed $N$ into its double. Let $p \in f^{-1}(\partial X), q = f(p) \in \partial X$ and choose a chart $\psi \colon V \to N$ around $q$, such that 
            \[
            \psi^{-1}(X) = V \cap \{(x_1,\ldots,x_n) \mid x_1 \leq 0, x_2 = \ldots = x_{s+1} = 0\},
            \]
            where $n = \mathrm{dim}(N)$ and $s = \mathrm{codim}(X)$. Since $f$ is a submersion we can find a chart $\phi \colon U \to M$ around $p$, such that 
            \[
            \psi^{-1} \circ f \circ \phi \colon U \to V
            \]
            is given by projection onto the first $n$ coordinates, i.e. $(x_1,\ldots,x_m) \mapsto (x_1,\ldots, x_n)$ for $m = \mathrm{dim}(M) \geq n$. Thus in the coordinates of $\phi$, $f^{-1}(X)$ is given by 
            \[
            \phi^{-1}(f^{-1}(X)) = U \cap \{(x_1,\ldots,x_m) \mid x_1 \leq 0, x_2 = \ldots = x_{s+1} = 0\},
            \]
            such that $\phi$ is a submanifold chart for $f^{-1}(X)$ around $p$ and $\partial(f^{-1}(X) \cap \phi(U)) = f^{-1}(X) \cap \phi(U)$. Similiarly we construct submanifold charts around points $p \in f^{-1}(\mathrm{int}(X))$ from which the statement follows.
        \end{proof}

        \begin{remark}
            The above Lemma becomes false if $\partial M \neq \emptyset$. Consider for example the inclusion $f \colon [0,1] \hookrightarrow \R$ and $S := [1,2] \subseteq \R$. Then 
            \[
            \partial f^{-1}(S) = \emptyset \neq f^{-1}(\partial S).
            \]
            Or even worse, consider the $2$-dimensional half space $\H^2 := \{(x,y) \in \R^2 \mid x \leq 0\}$ and the submersion $f \colon \H^2 \to \R, \; (x,y) \mapsto y$. Then $f^{-1}([0,1])$ is not even a manifold anymore.
            
            Moreover we cannot replace the sumbersion condition by only a transversality condition. Consider for example the map 
            \[
            f \colon \R \to \R^2, x \mapsto (0,x),
            \]
            which is transversal to $S := [0,1] \times \{0\} \subseteq \R^2$. Then again $\partial f^{-1}(S) = \emptyset \neq f^{-1}(\partial S)$.
        \end{remark}
    
 